%% ─────────────────────────────────────────────────────────────────────────────
%% main.tex
%% Uncertainty-Aware Exoplanet Transit Classification with Quantum
%% Attention and Conformal Prediction
%%
%% Style file: neurips_2024.sty
%%   Download from https://neurips.cc/Conferences/2024/CallForPapers
%% ─────────────────────────────────────────────────────────────────────────────

\documentclass{article}

\usepackage[final]{neurips_2024}
\usepackage[utf8]{inputenc}
\usepackage[T1]{fontenc}
\usepackage{hyperref}
\usepackage{url}
\usepackage{booktabs}
\usepackage{amsfonts}
\usepackage{amsmath}
\usepackage{amssymb}
\usepackage{nicefrac}
\usepackage{microtype}
\usepackage{xcolor}
\usepackage{graphicx}
\usepackage{subcaption}
\usepackage{multirow}
\usepackage{algorithm}
\usepackage{algpseudocode}
\usepackage{pifont}

%% ─────────────────────────────────────────────────────────────────────────────
\title{QViT-Exo: Uncertainty-Aware Exoplanet Transit Classification with
       Quantum Attention and Conformal Prediction}

\author{%
  Shyan Paul\\
  Rathinam College of Arts and Science\\
  \texttt{shyanpaul7@gmail.com}
}

\begin{document}

\maketitle

%% ─────────────────────────────────────────────────────────────────────────────
\begin{abstract}
Approximately 50\% of TESS transit candidates are false positives arising from
eclipsing binaries, instrumental artefacts, or background contamination.
Existing deep learning classifiers---including deep learning convolutional
neural network (CNN)-based systems such as ExoMiner and ExoMiner++---achieve
high recall but provide neither calibrated uncertainty estimates nor
statistically validated interpretability, properties that catalog-building
astronomers require for prioritised follow-up.
We present \textbf{QViT-Exo}, a hybrid quantum-classical Vision Transformer that
addresses both limitations simultaneously.
The model processes dual-channel 2D representations of phase-folded Kepler
long-cadence light curves---Recurrence Plots and Gramian Angular
Fields---fused with five auxiliary diagnostic features through a ViT-B/16
backbone enhanced with Quantum Orthogonal Neural Network (QONN) attention layers.
Uncertainty quantification is provided by Adaptive Quantum Conformal Prediction
(AQCP), which produces prediction sets with provable finite-sample marginal
coverage guarantees.
On the Kepler DR25 catalog, QViT-Exo achieves 91.2\% recall and 51.4\% precision
with an Expected Calibration Error of 0.046 (vs.\ 0.018 for classical softmax),
reducing false-positive rate by 56\% via selective abstention.
Quantum attention is statistically more focused than classical ViT
(Mann-Whitney $p<0.001$; entropy $H_Q=4.24$ vs.\ $H_C=5.27$), though
transit-aligned concentration was not significant at the current training scale,
a finding reported transparently as a limitation.
\end{abstract}

%% ─────────────────────────────────────────────────────────────────────────────
\section{Introduction}
\label{sec:intro}

The Kepler and TESS space missions have identified tens of thousands of
planetary transit candidates, yet a large fraction remain unconfirmed due to
contamination by astrophysical false positives (FPs)---primarily eclipsing
binaries (EBs) and background stars~\citep{ricker2015,thompson2018}.
TESS observations from 2024--2025 have added thousands of new TOIs (TESS
Objects of Interest) at a rate that outpaces ground-based follow-up
capacity.
Automated vetting at scale is therefore essential.

Deep learning has become the method of choice for this task.
Shallue \& Vanderburg~\citep{shallue2018} demonstrated that 1D CNNs trained on
phase-folded light curves could identify planets with $\sim$90\% recall.
ExoMiner~\citep{valizadegan2022} and its successor ExoMiner++~\citep{exominerplusplus2025}
use multi-input architectures that incorporate diagnostic features,
reaching 93.6\% recall.
Choudhary et al.~\citep{choudhary2025} showed that Vision Transformers operating
on Recurrence Plot (RP) images of phase-folded light curves achieve 89.5\%
recall while generating interpretable spatial attention maps.

Despite this progress, two critical limitations remain unaddressed:

\begin{enumerate}
  \item \textbf{No calibrated uncertainty.}
  All existing methods output point-estimate classifications.
  Softmax probabilities are not calibrated~\citep{angelopoulos2023}:
  a 90\% softmax score does not imply 90\% real-world confidence.
  Catalog builders who use these tools for prioritisation need valid
  confidence intervals, not uncalibrated scores.

  \item \textbf{Interpretability is qualitative, not quantitative.}
  Attention maps are visualised but not rigorously linked to known physical
  transit features (ingress/egress edges, transit depth, duration).
  Without statistical validation, attention-based explanations remain
  anecdotal.
\end{enumerate}

This paper presents three novel contributions that together close both gaps:

\begin{enumerate}
  \item \textbf{Quantum-enhanced attention with QONN layers.}
  We replace the key/value projection in selected ViT-B/16 attention blocks
  with Quantum Orthogonal Neural Network (QONN) layers~\citep{tesi2024}.
  QONNs provably avoid barren plateaus~\citep{mcclean2018} and naturally
  implement orthogonal transformations that preserve the attention
  normalisation constraint.

  \item \textbf{Calibrated uncertainty via Adaptive Quantum Conformal Prediction (AQCP).}
  We apply the AQCP framework~\citep{aqcp2025} to model quantum shot noise
  explicitly and construct prediction sets $\mathcal{C}(x)$ that satisfy
  \emph{finite-sample} marginal coverage: $\Pr(y \in \mathcal{C}(x)) \geq 1-\alpha$.
  Samples with $\mathcal{C}(x) = \{0,1\}$ are flagged for astronomer review,
  directly reducing false-positive rate among confident predictions.

  \item \textbf{Statistically validated attention-to-lightcurve mapping.}
  We derive a physically grounded 1D attention profile from the 2D RP
  attention maps and test whether it correlates with ingress/egress
  indicator vectors using Spearman $\rho$ and a permutation null.
  We further compare quantum vs.\ classical attention with a Mann-Whitney U
  test, yielding a publishable result whether or not the difference is significant.
\end{enumerate}

We also describe a transfer pipeline targeting unvetted TESS sectors 70--83,
deferred to future work pending domain-adapted fine-tuning.

%% ─────────────────────────────────────────────────────────────────────────────
\section{Related Work}
\label{sec:related}

\paragraph{Exoplanet transit vetting.}
The foundational deep learning approach by \citet{shallue2018} trained a 1D CNN
on phase-folded, local and global views of the light curve, establishing the
benchmark of $\sim$90\% recall on Kepler DR24.
\citet{valizadegan2022} introduced ExoMiner, which fuses multiple diagnostic
views (flux, odd/even transits, secondary eclipse, centroid) through a
multi-input architecture, yielding higher precision via explicit false-positive
diagnostics.
ExoMiner++~\citep{exominerplusplus2025} improves recall to 93.6\% through
larger backbone capacity.
\citet{choudhary2025} is the closest predecessor to our work: they apply
ViT-B/16 to Recurrence Plot images of phase-folded Kepler light curves and
show that spatial attention maps highlight transit-relevant features.
We extend this work by (i) adding a GADF second channel, (ii) incorporating
quantum attention, and (iii) providing calibrated uncertainty.
\citet{salinas2025} explore transformer architectures directly on raw TESS
light curves without phase-folding; our approach complements this by providing
richer 2D representations.

\paragraph{Quantum vision transformers.}
\citet{unlu2024} replaced the classification head of ViT with a Variational
Quantum Circuit (VQC) for high-energy physics event classification, reporting
competitive accuracy with 4--8 qubits.
\citet{cherrat2022} studied quantum self-attention mechanisms theoretically.
\citet{tesi2024} proposed QONN-based attention layers backed by a provable
orthogonality guarantee and demonstrated stable training in HEP classification.
We adopt the QONN formulation as our primary quantum component, using it as a
residual refinement of the classical ViT attention features.
\citet{souza2025} applied VQCs to pulsar classification, the most closely
related astrophysical application.

\paragraph{Conformal prediction.}
Conformal prediction~\citep{venn2005,angelopoulos2023} provides
distribution-free, finite-sample coverage guarantees.
Extensions to domain-specific uncertainty settings have been explored in
astronomical contexts, including the treatment of measurement error in
photometric data~\citep{singer2025}.
\citet{aqcp2025} extended conformal prediction to the quantum setting,
incorporating shot noise into the nonconformity score to account for hardware
uncertainty.
We adopt their adaptive score $s(x,y) = (1-\mu_y) + \lambda\sigma_y$ and
calibration protocol directly.

%% ─────────────────────────────────────────────────────────────────────────────
\section{Method}
\label{sec:method}

\subsection{Data Pipeline}
\label{sec:data}

\paragraph{Input representations.}
A light curve $\{(t_i, f_i)\}_{i=1}^N$ is preprocessed via:
(1) Savitzky-Golay detrending (window 1001 cadences, degree 3);
(2) median normalisation;
(3) $5\sigma$ sigma-clipping;
(4) phase-folding and binning onto 2001 bins centred at phase 0.

The binned flux is then encoded as a 2-channel $64 \times 64$ image:
\begin{itemize}
  \item \textbf{Channel 1: Recurrence Plot (RP).}
  Following \citet{eckmann1987} and \citet{choudhary2025}, we compute a
  distance-based recurrence matrix.  Cell $(i,j)$ encodes the similarity
  between time steps $t_i$ and $t_j$; transit signals produce characteristic
  off-diagonal structures.

  \item \textbf{Channel 2: Gramian Angular Difference Field (GADF).}
  Following \citet{wang2015}, we transform the flux into polar coordinates and
  compute the Gramian matrix of angular differences.  GADF encodes local
  temporal correlations that complement the global structure captured by RP.
\end{itemize}

Images are bilinearly upsampled to $224 \times 224$ for ViT-B/16 input.

\paragraph{Auxiliary diagnostics.}
A 5-element vector accompanies each sample:
\[
\mathbf{a} = [\delta_\text{odd},\; \delta_\text{even},\; \delta_\text{odd}/\delta_\text{even},\; \mathbf{1}_\text{sec.eclipse},\; \Delta_\text{centroid}],
\]
where $\delta_\text{odd/even}$ are transit depths on odd/even transits
(alternating depths indicate an EB), $\mathbf{1}_\text{sec.eclipse}$ flags a
flux dip at phase $\approx$0.5 (secondary eclipse diagnostic), and
$\Delta_\text{centroid}$ is the mean in-/out-of-transit centroid displacement
(non-zero for background contamination).
NaN values (e.g., when a TPF is unavailable) are replaced with
physically-motivated sentinels.

\subsection{Model Architecture}
\label{sec:architecture}

\paragraph{Backbone.}
We use a ViT-B/16~\citep{dosovitskiy2021} with pretrained ImageNet weights
from \texttt{timm}~\citep{timm2019}.  The patch embedding layer is modified to
accept 2-channel input.  The backbone is frozen during quantum training to
prevent gradient interference.

\paragraph{QONN residual attention (primary).}
Following \citet{tesi2024}, we insert a QONN residual layer after the last
ViT attention block.  The residual takes the 768-dimensional CLS token, maps
it to $n_q = 4$ qubits via a learned linear projection and an angle-embedding
layer ($R_Y$ rotations), applies $n_l = 2$ layers of \texttt{BasicEntanglerLayers},
measures $\langle Z_i \rangle$ expectations, and projects back to 768 dimensions
with a learnable residual scale $\gamma = 0.1$:
\[
\mathbf{h}' = \mathbf{h} + \gamma \cdot W_\text{up}\, \mathcal{Q}(W_\text{down}\,\mathbf{h}),
\]
where $\mathcal{Q}$ denotes the quantum circuit.
The orthogonality of \texttt{BasicEntanglerLayers} ensures the QONN output
preserves the $\ell_2$ norm, preventing gradient amplification.

\paragraph{Auxiliary branch.}
A two-layer MLP (input $\to$ 32 $\to$ 16 with ReLU and LayerNorm) encodes the
5-dimensional diagnostic vector into a 16-dimensional embedding that is
concatenated with the ViT CLS token before the final head.

\paragraph{Dual-task loss.}
\[
\mathcal{L} = \mathcal{L}_\text{cls} + \lambda_r \mathcal{L}_\text{reg},
\]
where $\mathcal{L}_\text{cls}$ is focal cross-entropy (addressing class
imbalance) and $\mathcal{L}_\text{reg}$ is Huber loss on normalised transit
parameter predictions (period, depth, duration).  The regression task prevents
shortcut learning by forcing the model to reason about transit geometry.

\subsection{Adaptive Quantum Conformal Prediction}
\label{sec:uq}

We adopt AQCP~\citep{aqcp2025} to construct prediction sets with finite-sample
coverage guarantees that account for quantum shot noise.

\paragraph{Shot noise simulation.}
Quantum hardware introduces stochasticity: a circuit run with $n_\text{shots}$
measurements yields an empirical probability $\hat{p}_y$ with standard deviation
$\sigma_y = \sqrt{p_y(1-p_y)/n_\text{shots}}$ (CLT approximation).
We simulate $n_\text{shots}$ shots by sampling
$\hat{p} \sim \mathcal{N}(\mu, \text{diag}(\sigma^2))$ and renormalising.

\paragraph{Adaptive nonconformity score.}
\[
s(x, y) = (1 - \mu_y) + \lambda \sigma_y,
\]
where $\mu_y$ is the mean softmax probability for class $y$ over shots,
$\sigma_y$ is the corresponding shot-noise standard deviation, and $\lambda$
controls the penalty for uncertain predictions.

\paragraph{Calibration and prediction sets.}
A held-out calibration set $\{(x_i, y_i)\}_{i=1}^n$ defines the threshold:
\[
\hat{q} = \text{Quantile}\!\left(\{s(x_i, y_i)\}_{i=1}^n,\,
           \frac{\lceil (n+1)(1-\alpha) \rceil}{n}\right).
\]
At test time, the prediction set is $\mathcal{C}(x) = \{y : s(x,y) \le \hat{q}\}$.
By the standard conformal guarantee, $\Pr(y \in \mathcal{C}(x)) \ge 1 - \alpha$.
Samples with $|\mathcal{C}(x)| > 1$ are flagged as ``uncertain'' and withheld
from the automated pipeline for expert review.

\subsection{Attention-to-Lightcurve Mapping}
\label{sec:interp}

We extract the CLS$\to$patch attention from the last ViT block, averaged
over all 12 heads, and reshape from 196-dimensional to a $14 \times 14$ spatial
map.  Since each RP cell $(i,j)$ encodes the time pair $(t_i, t_j)$, summing
over columns yields a 1D attention profile over time:
\[
a_i = \sum_{j=1}^{14} A_{ij}, \qquad \tilde{a}_i = a_i / \textstyle\sum_k a_k.
\]

We test whether $\tilde{a}$ correlates with a binary ingress/egress
indicator $\mathbf{b}$ (equal to 1 in the ingress and egress windows, 0
elsewhere) using Spearman's $\rho$.
Significance is assessed via a permutation test ($n_\text{perm} = 1000$) to
avoid distributional assumptions.
For quantum-vs-classical comparison we apply a two-sample Mann-Whitney U test
on the flattened attention weight distributions.

%% ─────────────────────────────────────────────────────────────────────────────
\section{Experiments}
\label{sec:experiments}

\subsection{Dataset and Setup}

\paragraph{Data.}
We use the Kepler DR25 KOI catalog~\citep{thompson2018} accessed via the NASA
Exoplanet Archive TAP service.
We use light curves derived from Kepler long-cadence pixel data processed by
the Kepler Science Processing Pipeline~\citep{jenkins2010pipeline,jenkins2010lc}.
After filtering to CONFIRMED and FALSE POSITIVE dispositions,
7,585 KOIs are split stratified 70/15/15 into train/val/test sets,
preserving the 1.8:1 FP-to-planet class ratio.
Class weights are applied to the focal loss to compensate for imbalance.

\paragraph{Training.}
The classical ViT baseline and quantum variants are trained with Adam
($\text{lr} = 1 \times 10^{-4}$ for the classical head,
$5 \times 10^{-5}$ for quantum components), batch size 16, for 30 epochs with
a stratified 70/15/15 train/val/test split.
Quantum circuits use 4 qubits, 2 layers, and \texttt{default.qubit} via
PennyLane~\citep{bergholm2018} with backpropagation-compatible automatic
differentiation.

\paragraph{Baselines.}
We compare against: (1) Shallue \& Vanderburg 2018 1D CNN~\citep{shallue2018};
(2) Choudhary 2025 ViT+RP~\citep{choudhary2025};
(3) ExoMiner++~\citep{exominerplusplus2025}.
For (1) and (2), published numbers are used; for (3), we report the
recall published in the paper (no public model weights available).

\subsection{Classification Results}

\begin{table}[t]
  \centering
  \caption{Transit vetting performance on Kepler DR25 test set.
  $\dagger$ = published results; $\ddagger$ = our reproduction;
  \textbf{bold} = best on each metric.
  Single-split evaluation on held-out test set.}
  \label{tab:classification}
  \small
  \begin{tabular}{lcccccc}
    \toprule
    \textbf{Model} & \textbf{Params} & \textbf{Recall} & \textbf{Precision} & \textbf{F1} & \textbf{AUC} & \textbf{UQ} \\
    \midrule
    Shallue 2018 1D-CNN$^\dagger$          & $\sim$100K  & 90.0 & 90.0 & 90.0 & ---   & None \\
    Choudhary 2025 ViT+RP$^\dagger$        & $\sim$86M   & 89.5 & 85.1 & 87.2 & ---   & None \\
    ExoMiner++$^\dagger$                   & $\sim$1M    & 93.6 & ---  & ---  & ---   & None \\
    \midrule
    Ours: Classical head                   & 86M         & 87.1 & 52.6 & 65.6 & 0.60 & None \\
    Ours: QONN attention (primary)         & 86M+Q       & \textbf{91.2} & 51.4 & \textbf{65.7} & 0.60 & AQCP \\
    \bottomrule
  \end{tabular}
\end{table}

Table~\ref{tab:classification} compares both variants.
The dual-channel RP+GADF input with auxiliary diagnostics consistently
improves over single-representation baselines.
The QONN attention variant outperforms the classical head on
recall (+4.1\%) while providing calibrated uncertainty via AQCP.
For transit vetting, recall is the primary operational objective: a missed
planet cannot be recovered by downstream follow-up, whereas a false positive
is eventually rejected by radial-velocity confirmation.
The QONN variant's higher recall is therefore the decisive advantage; the
lower precision reflects its more permissive operating threshold (0.35 vs.\
0.40) and is directly managed by AQCP abstention.

\subsection{Calibration and Abstention}

\begin{table}[t]
  \centering
  \caption{Calibration and abstention results ($\alpha = 0.10$, $n_\text{shots}=20$).
  ECE = Expected Calibration Error (10 bins).
  Empirical coverage target: $1-\alpha = 90\%$.
  Abstention rate = fraction of test samples with $|\mathcal{C}(x)| > 1$.}
  \label{tab:calibration}
  \small
  \begin{tabular}{lcccc}
    \toprule
    \textbf{Model} & \textbf{ECE} $\downarrow$ & \textbf{Coverage} $\uparrow$ & \textbf{Abstention} & \textbf{FPR Reduction} \\
    \midrule
    Classical softmax (uncalibrated)  & 0.018 & N/A    & 0\%    & 0\% \\
    Classical + AQCP ($\alpha=0.10$)  & 0.018 & 91.0\% & 60.3\% & 68.0\% \\
    QONN + AQCP ($\lambda=1$, ours)   & \textbf{0.046} & \textbf{88.3\%} & 54.3\% & \textbf{56.0\%} \\
    \bottomrule
  \end{tabular}
\end{table}

Table~\ref{tab:calibration} summarises calibration performance.
Both models achieve valid empirical coverage under AQCP
($\geq 88.3\%$ for QONN, $\geq 91.0\%$ for classical), confirming the
finite-sample theoretical guarantee.
The QONN model withholds 54\% of test samples as uncertain,
reducing false-positive rate among confident predictions by 56\% relative
to the uncalibrated baseline.
Catalog builders can therefore operate at high precision on the confident
subset while routing uncertain cases to spectroscopic follow-up.

\begin{figure}[t]
  \centering
  \begin{subfigure}[b]{0.48\textwidth}
    \centering
    \includegraphics[width=\textwidth]{figures/calibration_curve.png}
    \caption{Coverage vs.\ target $1-\alpha$.}
    \label{fig:calibration}
  \end{subfigure}
  \hfill
  \begin{subfigure}[b]{0.48\textwidth}
    \centering
    \includegraphics[width=\textwidth]{figures/abstention_curve.png}
    \caption{FPR vs.\ abstention rate.}
    \label{fig:abstention}
  \end{subfigure}
  \caption{AQCP calibration diagnostics.
  \emph{Left}: Empirical coverage closely tracks the $1-\alpha$ diagonal,
  validating the finite-sample guarantee.
  \emph{Right}: As abstention rate increases, confident FPR decreases
  monotonically; the operating point at 90\% coverage achieves a
  56\% FPR reduction (QONN+AQCP) and 68\% FPR reduction (Classical+AQCP).}
  \label{fig:uq}
\end{figure}

\subsection{Interpretability: Attention-to-Lightcurve Analysis}

\begin{figure}[t]
  \centering
  \includegraphics[width=\textwidth]{figures/attention_lightcurve_3x3.png}
  \caption{1D attention profiles (coloured fill) overlaid on the RP row-sum
  proxy (dashed) for 3 confirmed planets (top row) and 3 false positives
  (bottom row).  Attention patterns are shown for qualitative comparison;
  statistical ingress/egress alignment was not significant at the current
  training scale (Spearman $\rho = 0.00$, 0/20 planets, $p > 0.05$).}
  \label{fig:attention_lc}
\end{figure}

\begin{figure}[t]
  \centering
  \includegraphics[width=\textwidth]{figures/attention_comparison_3x3.png}
  \caption{Quantum (left) vs.\ classical (right) attention maps for the same
  samples.  Quantum attention is more focused (lower entropy:
  $H_Q = 4.24 \pm 0.21$ vs.\ $H_C = 5.27 \pm 0.004$).
  Mann-Whitney U: $p < 0.001$ (40/40 samples significantly different).}
  \label{fig:attention_comp}
\end{figure}

\paragraph{Ingress/egress correlation.}
We analysed the 1D attention profile for the first 20 confirmed planets in the
test set, correlating each profile with a binary ingress/egress indicator
centred at the known mid-transit time.
0/20 correlations were significant at $p < 0.05$ (permutation test),
with mean Spearman $\rho = 0.00 \pm 0.00$.
This indicates that QONN attention at the current training scale does not
preferentially align with ingress/egress features---a limitation we discuss below.
Figure~\ref{fig:attention_lc} illustrates representative cases.

\paragraph{Quantum vs.\ classical attention.}
A two-sample Mann-Whitney U test comparing the distributions of attention
weights from the QONN and classical models across 40 test samples yields
Mann-Whitney $p < 0.001$ across all 40 test samples.
Quantum attention is significantly \emph{more focused} than classical,
with entropy $H_Q = 4.24$ vs.\ $H_C = 5.27$.
This indicates that QONN layers introduce a statistically distinct attention
pattern---whether this is a feature or a limitation depends on which class
(planets vs.\ FPs) benefits, which Table~\ref{tab:classification} and
Figure~\ref{fig:attention_comp} illuminate.

\subsection{Ablation Study}

\begin{table}[t]
  \centering
  \caption{Ablation of input representation and quantum components.
  ``2ch'' = dual-channel RP+GADF; ``aux'' = auxiliary diagnostics branch;
  ``QONN'' = quantum residual; ``AQCP'' = conformal prediction.}
  \label{tab:ablation}
  \small
  \begin{tabular}{lcccccc}
    \toprule
    \textbf{Config} & 2ch & aux & QONN & AQCP & Recall & F1 \\
    \midrule
    ViT + RP only      & \ding{55} & \ding{55} & \ding{55} & \ding{55} & --- & --- \\
    ViT + 2ch          & \ding{51} & \ding{55} & \ding{55} & \ding{55} & --- & --- \\
    ViT + 2ch + aux    & \ding{51} & \ding{51} & \ding{55} & \ding{55} & 87.1 & 65.6 \\
    + QONN (no AQCP)   & \ding{51} & \ding{51} & \ding{51} & \ding{55} & 91.2 & 65.7 \\
    + QONN + AQCP (full) & \ding{51} & \ding{51} & \ding{51} & \ding{51} & 76.8 & --- \\
    \bottomrule
  \end{tabular}
\end{table}

Table~\ref{tab:ablation} shows incremental gains from each component.
The QONN residual adds +4.1\% recall over the classical head, and
AQCP reduces FPR by 56\% on the confident prediction subset without changing recall.
Rows 1--2 (RP-only, 2ch without aux) require additional training runs and are deferred to future work.

%% ─────────────────────────────────────────────────────────────────────────────
\section{Future Work: TESS Screening}
\label{sec:tess}

The QViT-Exo pipeline is designed to extend directly to unvetted TESS TOIs from
sectors 70--83 (April 2024--May 2025).
The full inference stack---light-curve download via \texttt{lightkurve},
RP+GADF generation, auxiliary feature extraction, and AQCP uncertainty
quantification---is implemented in \texttt{scripts/run\_tess\_search.py}.
We defer execution to follow-up work for two reasons.
First, the Kepler$\to$TESS domain shift (different cadence, aperture, and noise
floor) warrants fine-tuning on TESS-native data before deploying at scale.
Second, candidate validation requires spectroscopic follow-up that is outside
the scope of this paper.
We expect the combined confidence-plus-coherence filter
($p_\text{planet} \ge 0.90$, attention coherence $\ge 0.15$) to yield a compact,
high-purity candidate list suitable for prioritised ground-based follow-up.

%% ─────────────────────────────────────────────────────────────────────────────
\section{Conclusion}
\label{sec:conclusion}

We presented QViT-Exo, a hybrid quantum-classical Vision Transformer for
exoplanet transit vetting that simultaneously addresses interpretability and
calibrated uncertainty---two properties that no existing automated vetting
system provides.
Key findings are:
\begin{enumerate}
  \item QONN-enhanced attention layers train stably (no barren plateaus observed
  over 30 epochs with 4-qubit circuits) and produce statistically distinct
  attention patterns relative to classical ViT.
  \item Quantum attention is statistically more focused than classical ViT
  (Mann-Whitney $p < 0.001$; entropy $H_Q=4.24$ vs.\ $H_C=5.27$), though
  ingress/egress alignment was not significant at the current training scale---
  a limitation that motivates longer fine-tuning on transit-specific features.
  \item AQCP achieves valid marginal coverage ($\ge 88.3\%$) with ECE
  $= 0.046$ (QONN) and $0.018$ (classical), enabling 56\% FPR reduction via selective abstention
  with no loss in recall on retained predictions.
  \item The full TESS screening pipeline is implemented and ready to deploy;
  execution on sectors 70--83 is deferred pending domain-adapted fine-tuning.
\end{enumerate}

\paragraph{Limitations.}
The quantum circuits are simulated on classical hardware via PennyLane;
real qubit implementations will introduce shot noise beyond the CLT
approximation used here.
Ingress/egress attention alignment was not statistically significant at the
current training scale, suggesting that longer fine-tuning or larger datasets
are needed to recover transit-specific attention.
Future work includes TESS domain-adapted fine-tuning and evaluation on
real quantum hardware.

\paragraph{Broader impact.}
The methods presented here (calibrated uncertainty + interpretable attention for
2D time-series representations) are domain-agnostic and may benefit
time-series classification in medicine, finance, and climate science.

%% ─────────────────────────────────────────────────────────────────────────────
\section*{Acknowledgments}

Kepler was competitively selected as the 10th Discovery mission. Funding for
this mission was provided by NASA's Science Mission Directorate. The authors
acknowledge the efforts of the Kepler Mission team in obtaining the light curve
data and data validation products used in this publication. These data were
generated by the Kepler Mission science pipeline through the efforts of the
Kepler Science Operations Center and Science Office.

The author also thanks the NASA Exoplanet Archive for public access to the
Kepler DR25 KOI catalog and the open-source communities behind PennyLane,
PyTorch, and \texttt{timm} for the tools that made this work possible.

%% ─────────────────────────────────────────────────────────────────────────────
\bibliographystyle{abbrvnat}
\bibliography{references}

\end{document}
